% =============================================================
% resume.tex — Résumé en français
% Équivalent strict de frontmatter/abstract.tex : toute modification
% ici doit être répercutée à l'identique dans l'abstract anglais.
% Le résumé et l'abstract partagent une seule page (norme ESPRIT §H,
% « Résumé / Abstract : 1 page ») : abstract.tex enchaîne sur cette page
% sans \chapter*, avec un titre de même graisse et de même corps.
% =============================================================
\chapter*{Résumé}
\addcontentsline{toc}{chapter}{Résumé}
\thispagestyle{plain}
\vspace*{-10pt}

{\footnotesize

Ce mémoire présente la conception, la réalisation et l'évaluation d'un socle d'hébergement
sécurisé sur Google Cloud pour MENAL-SARL. Le point de départ est une application pilote sur un
serveur unique, sans reproductibilité complète ni chaîne de détection. La démarche associe audit
technique, analyse de risque, vingt exigences vérifiables, architecture Zero Trust,
infrastructure décrite en code et portes de contrôle DevSecOps bloquantes. Le socle obtenu
sépare les identités de service, déploie sans clé permanente et exécute sept règles de détection
sur une collecte filtrée de journaux ; un modèle spécialisé rapproche les alertes des techniques
MITRE ATT\&CK, hors du score d'incident. Le 19/08/2026, treize requêtes malveillantes d'un
scénario de bout en bout ont toutes été bloquées en HTTP~403, retrouvées dans les journaux et
agrégées par la règle R2 quinze minutes plus tard. Le 24/08/2026, une sortie réseau non
autorisée et une connexion croisée entre locataires ont été refusées et journalisées, les deux
horodatages se recoupant ; le contrôle d'isolation des données du locataire a réussi six
vérifications sur six. La restauration de la base a été mesurée à 32~min~45~s, sans perte de
données. Les limites principales sont la sortie réseau non gouvernée de deux charges de travail
détachées du sous-réseau, un motif de traversée de chemin encore partiellement filtré, l'instance
de base de données encore partagée entre locataires et l'absence de mesures complètes de coût et
de qualité sémantique. Le socle est reproductible et exploitable en recette ; l'apport principal
n'est pas un composant nouveau mais une méthode : chaque contrôle est associé au test qui en
démontre l'efficacité, et le mémoire publie les contrôles qui tiennent comme ceux qui ne
tiennent pas encore.

\par\vspace{4pt}
\noindent\textbf{Mots-clés :} Zero Trust, DevSecOps, Google Cloud, infrastructure as code,
SIEM, MITRE ATT\&CK, sécurité multi-locataire.

}
